% ============================================================
\chapter{Espaces Probabilis\'es}
\label{ch:espaces}
% ============================================================

Le calcul des probabilit\'es est n\'e dans les tripots du XVII\ieme{} si\`ecle.
Blaise Pascal et Pierre de Fermat, dans leur c\'el\`ebre correspondance de
1654, cherchaient \`a r\'esoudre le \og probl\`eme des partis\fg~: comment
r\'epartir \'equitablement la mise d'un jeu interrompu~? Mais il a fallu
pr\`es de trois si\`ecles pour que la th\'eorie des probabilit\'es acqui\`ere des
fondations math\'ematiques rigoureuses. C'est Andre\"\i{} Kolmogorov qui, en
1933, dans ses \emph{Grundbegriffe der Wahrscheinlichkeitsrechnung}, posa
les axiomes d\'efinitifs~: une probabilit\'e n'est rien d'autre qu'une mesure
de masse totale $1$ sur un espace mesurable.

Ce premier chapitre construit ces fondations pierre par pierre~: l'univers
$\Omega$ des possibles, la tribu $\mathcal{F}$ des \'ev\'enements
observables, et la mesure de probabilit\'e $\PP$ qui attribue \`a chaque
\'ev\'enement un degr\'e de vraisemblance.

\begin{intuition}
Ce premier chapitre pose les fondations math\'ematiques de la th\'eorie des probabilit\'es.
Nous y d\'efinissons la notion d'espace probabilis\'e $(\Omega, \mathcal{F}, \PP)$,
pr\'esentons les axiomes de Kolmogorov et \'etablissons les premi\`eres propri\'et\'es
de la mesure de probabilit\'e.
\end{intuition}

% ------------------------------------------------------------
\section{Exp\'erience al\'eatoire et univers}
% ------------------------------------------------------------

Avant toute formalisation, posons la question na\"ive~: qu'est-ce qu'une
exp\'erience al\'eatoire~? C'est une exp\'erience dont nous ne pouvons pas
pr\'edire le r\'esultat, mais dont nous pouvons d\'ecrire tous les r\'esultats
possibles. Un lancer de d\'e, le cours d'une action boursi\`ere demain, la
dur\'ee de vie d'un transistor --- autant de situations o\`u l'incertitude
est irr\'eductible, mais o\`u la structure de l'incertitude est accessible.

\begin{definition}[Exp\'erience al\'eatoire]
Une \textbf{exp\'erience al\'eatoire} est une exp\'erience dont le r\'esultat ne peut
\^etre pr\'edit avec certitude avant sa r\'ealisation, mais dont on peut d\'ecrire
l'ensemble de tous les r\'esultats possibles.
\end{definition}

\begin{definition}[Univers]
L'\textbf{univers} (ou \textbf{espace fondamental}) est l'ensemble $\Omega$ de tous
les r\'esultats possibles d'une exp\'erience al\'eatoire. Chaque \'el\'ement
$\omega \in \Omega$ est appel\'e un \textbf{r\'esultat \'el\'ementaire} ou
\textbf{\'epreuve}.
\end{definition}

\begin{exemple}
\begin{enumerate}
    \item \textbf{Lancer d'un d\'e~:} $\Omega = \{1, 2, 3, 4, 5, 6\}$.
    \item \textbf{Lancer d'une pi\`ece~:} $\Omega = \{\text{Pile}, \text{Face}\}$.
    \item \textbf{Dur\'ee de vie d'un composant~:} $\Omega = [0, +\infty)$.
    \item \textbf{Suite infinie de lancers de pi\`ece~:}
          $\Omega = \{0, 1\}^{\N} = \{(\omega_1, \omega_2, \ldots) : \omega_i \in \{0,1\}\}$.
\end{enumerate}
\end{exemple}

\begin{remarque}
L'univers $\Omega$ peut \^etre fini, d\'enombrable ou ind\'enombrable. Le choix de
$\Omega$ d\'epend de la mod\'elisation~: il n'est pas unique pour une exp\'erience donn\'ee.
\end{remarque}

% ------------------------------------------------------------
\section{Tribus et \'ev\'enements}
% ------------------------------------------------------------

Lorsque $\Omega$ est fini ou d\'enombrable, on peut souvent consid\'erer tous les
sous-ensembles de $\Omega$ comme des \'ev\'enements. Pour des espaces ind\'enombrables
(comme $\R$), il est n\'ecessaire de restreindre la classe des ensembles
\og mesurables \fg.

\begin{definition}[Tribu / $\sigma$-alg\`ebre]
\label{def:tribu}
Une \textbf{tribu} (ou $\sigma$-alg\`ebre) sur $\Omega$ est une collection
$\mathcal{F} \subseteq \mathcal{P}(\Omega)$ v\'erifiant~:
\begin{enumerate}[label=(\roman*)]
    \item $\Omega \in \mathcal{F}$\,;
    \item si $A \in \mathcal{F}$, alors $A^c \in \mathcal{F}$ \quad
          (stabilit\'e par compl\'ementation)\,;
    \item si $(A_n)_{n \geq 1}$ est une suite d'\'el\'ements de $\mathcal{F}$, alors
          $\displaystyle\bigcup_{n=1}^{\infty} A_n \in \mathcal{F}$ \quad
          (stabilit\'e par union d\'enombrable).
\end{enumerate}
\end{definition}

\begin{proposition}[Propri\'et\'es imm\'ediates]
Soit $\mathcal{F}$ une tribu sur $\Omega$. Alors~:
\begin{enumerate}[label=(\alph*)]
    \item $\emptyset \in \mathcal{F}$\,;
    \item $\mathcal{F}$ est stable par intersection d\'enombrable~:
          si $A_n \in \mathcal{F}$ pour tout $n$, alors
          $\bigcap_{n=1}^{\infty} A_n \in \mathcal{F}$\,;
    \item $\mathcal{F}$ est stable par diff\'erence~: si $A, B \in \mathcal{F}$,
          alors $A \setminus B \in \mathcal{F}$.
\end{enumerate}
\end{proposition}

\begin{proof}
(a) $\emptyset = \Omega^c \in \mathcal{F}$ par (i) et (ii).
(b) Par les lois de De Morgan, $\bigcap A_n = \left(\bigcup A_n^c\right)^c$,
ce qui est dans $\mathcal{F}$ par (ii) et (iii).
(c) $A \setminus B = A \cap B^c \in \mathcal{F}$ par (ii) et (b).
\end{proof}

\begin{definition}[\'Ev\'enement]
Un \textbf{\'ev\'enement} est un \'el\'ement de la tribu $\mathcal{F}$, c'est-\`a-dire
un sous-ensemble de $\Omega$ auquel on peut attribuer une probabilit\'e. On dit que
l'\'ev\'enement $A$ \textbf{se r\'ealise} si le r\'esultat $\omega$ de l'exp\'erience
appartient \`a $A$.
\end{definition}

\begin{exemple}[Tribus fondamentales]
\begin{enumerate}
    \item \textbf{Tribu grossi\`ere~:} $\mathcal{F} = \{\emptyset, \Omega\}$.
          C'est la plus petite tribu sur $\Omega$.
    \item \textbf{Tribu discr\`ete~:} $\mathcal{F} = \mathcal{P}(\Omega)$.
          C'est la plus grande tribu sur $\Omega$.
    \item \textbf{Tribu engendr\'ee~:} si $\mathcal{C}$ est une collection de
          sous-ensembles de $\Omega$, la tribu engendr\'ee $\sigma(\mathcal{C})$
          est la plus petite tribu contenant $\mathcal{C}$.
\end{enumerate}
\end{exemple}

\begin{definition}[Tribu bor\'elienne]
La \textbf{tribu bor\'elienne} de $\R$, not\'ee $\mathcal{B}(\R)$, est la tribu
engendr\'ee par les ouverts de $\R$~:
\[
    \mathcal{B}(\R) = \sigma\bigl(\{O \subseteq \R : O \text{ ouvert}\}\bigr).
\]
On montre que $\mathcal{B}(\R) = \sigma\bigl(\{]-\infty, a] : a \in \R\}\bigr)
= \sigma\bigl(\{]a, b[ : a < b\}\bigr)$.
\end{definition}

% ------------------------------------------------------------
\section{Axiomes de Kolmogorov}
% ------------------------------------------------------------

\begin{definition}[Mesure de probabilit\'e]
\label{def:proba}
Soit $(\Omega, \mathcal{F})$ un espace mesurable. Une \textbf{mesure de probabilit\'e}
est une application $\PP : \mathcal{F} \to [0,1]$ v\'erifiant les
\textbf{axiomes de Kolmogorov}~:
\begin{enumerate}[label=\bfseries (K\arabic*)]
    \item \textbf{Non-n\'egativit\'e~:} pour tout $A \in \mathcal{F}$,
          $\PP(A) \geq 0$.
    \item \textbf{Normalisation~:} $\PP(\Omega) = 1$.
    \item \textbf{$\sigma$-additivit\'e~:} pour toute suite $(A_n)_{n \geq 1}$
          d'\'ev\'enements \emph{mutuellement disjoints} (i.e. $A_i \cap A_j = \emptyset$
          pour $i \neq j$)~:
          \[
              \PP\!\left(\bigcup_{n=1}^{\infty} A_n\right)
              = \sum_{n=1}^{\infty} \PP(A_n).
          \]
\end{enumerate}
\end{definition}

\begin{definition}[Espace probabilis\'e]
Le triplet $(\Omega, \mathcal{F}, \PP)$ est appel\'e \textbf{espace probabilis\'e}.
\end{definition}

\begin{attention}[title=\textbf{Additivit\'e finie vs.\ $\sigma$-additivit\'e}]
L'additivit\'e finie ($\PP(A \cup B) = \PP(A) + \PP(B)$ pour $A \cap B = \emptyset$)
d\'ecoule de (K3). L'axiome de $\sigma$-additivit\'e est plus fort~: il garantit la
coh\'erence du passage \`a la limite et est essentiel pour les th\'eor\`emes limites
(chapitres~9--11).
\end{attention}

% ------------------------------------------------------------
\section{Propri\'et\'es de la mesure de probabilit\'e}
% ------------------------------------------------------------

\begin{theoreme}[Propri\'et\'es fondamentales]
\label{thm:prop_proba}
Soit $(\Omega, \mathcal{F}, \PP)$ un espace probabilis\'e. Pour tous
$A, B \in \mathcal{F}$~:
\begin{enumerate}[label=(\roman*)]
    \item $\PP(\emptyset) = 0$.
    \item $\PP(A^c) = 1 - \PP(A)$.
    \item Si $A \subseteq B$, alors $\PP(A) \leq \PP(B)$ \quad (monotonie).
    \item $\PP(A \cup B) = \PP(A) + \PP(B) - \PP(A \cap B)$.
    \item $\PP(A) \leq 1$.
\end{enumerate}
\end{theoreme}

\begin{proof}
\begin{enumerate}[label=(\roman*)]
    \item Posons $A_1 = \Omega$ et $A_n = \emptyset$ pour $n \geq 2$. La suite est
    disjointe, et $\bigcup A_n = \Omega$. Par (K3)~:
    $1 = \PP(\Omega) = \PP(\Omega) + \sum_{n \geq 2} \PP(\emptyset)$.
    Donc $\PP(\emptyset) = 0$.

    \item $\Omega = A \cup A^c$ (union disjointe), donc
    $1 = \PP(A) + \PP(A^c)$.

    \item $B = A \cup (B \setminus A)$ (union disjointe), donc
    $\PP(B) = \PP(A) + \PP(B \setminus A) \geq \PP(A)$.

    \item $A \cup B = A \cup (B \setminus A)$ et
    $B = (A \cap B) \cup (B \setminus A)$, d'o\`u
    $\PP(B \setminus A) = \PP(B) - \PP(A \cap B)$ et le r\'esultat suit.

    \item Cons\'equence de (ii)~: $\PP(A) = 1 - \PP(A^c) \leq 1$.
\end{enumerate}
\end{proof}

\begin{theoreme}[Formule d'inclusion-exclusion]
\label{thm:inclusion_exclusion}
Pour tous \'ev\'enements $A_1, \ldots, A_n \in \mathcal{F}$~:
\[
    \PP\!\left(\bigcup_{i=1}^{n} A_i\right)
    = \sum_{i} \PP(A_i)
    - \sum_{i < j} \PP(A_i \cap A_j)
    + \sum_{i < j < k} \PP(A_i \cap A_j \cap A_k)
    - \cdots
    + (-1)^{n+1}\, \PP(A_1 \cap \cdots \cap A_n).
\]
\end{theoreme}

\begin{corollaire}[In\'egalit\'e de Boole / sous-additivit\'e]
Pour toute suite $(A_n)_{n \geq 1}$ d'\'ev\'enements~:
\[
    \PP\!\left(\bigcup_{n=1}^{\infty} A_n\right)
    \leq \sum_{n=1}^{\infty} \PP(A_n).
\]
\end{corollaire}

% ------------------------------------------------------------
\section{Continuit\'e de la mesure de probabilit\'e}
% ------------------------------------------------------------

\begin{theoreme}[Continuit\'e monotone]
\label{thm:continuite}
Soit $(\Omega, \mathcal{F}, \PP)$ un espace probabilis\'e.
\begin{enumerate}[label=(\roman*)]
    \item \textbf{Continuit\'e croissante~:} si $A_1 \subseteq A_2 \subseteq \cdots$
    (suite croissante), alors
    \[
        \PP\!\left(\bigcup_{n=1}^{\infty} A_n\right)
        = \lim_{n \to \infty} \PP(A_n).
    \]
    \item \textbf{Continuit\'e d\'ecroissante~:} si $A_1 \supseteq A_2 \supseteq \cdots$
    (suite d\'ecroissante), alors
    \[
        \PP\!\left(\bigcap_{n=1}^{\infty} A_n\right)
        = \lim_{n \to \infty} \PP(A_n).
    \]
\end{enumerate}
\end{theoreme}

\begin{proof}
(i) Posons $B_1 = A_1$ et $B_n = A_n \setminus A_{n-1}$ pour $n \geq 2$. Les $B_n$
sont mutuellement disjoints et $\bigcup_{n=1}^{N} B_n = A_N$. Par $\sigma$-additivit\'e~:
\[
    \PP\!\left(\bigcup_{n=1}^{\infty} A_n\right)
    = \PP\!\left(\bigcup_{n=1}^{\infty} B_n\right)
    = \sum_{n=1}^{\infty} \PP(B_n)
    = \lim_{N \to \infty} \sum_{n=1}^{N} \PP(B_n)
    = \lim_{N \to \infty} \PP(A_N).
\]

(ii) Appliquer (i) \`a la suite croissante $A_1^c \subseteq A_2^c \subseteq \cdots$ et
utiliser $\PP(A_n^c) = 1 - \PP(A_n)$.
\end{proof}

% ------------------------------------------------------------
\section{Construction de probabilit\'es}
% ------------------------------------------------------------

\subsection{Probabilit\'e uniforme sur un ensemble fini}

\begin{definition}[Probabilit\'e uniforme]
Si $\Omega$ est un ensemble fini, la \textbf{probabilit\'e uniforme} (ou
\'equiprobabilit\'e) est d\'efinie par~:
\[
    \PP(A) = \frac{\abs{A}}{\abs{\Omega}} = \frac{\text{nombre de cas favorables}}%
    {\text{nombre de cas possibles}}.
\]
\end{definition}

\begin{intuition}[title=\textbf{Quand utiliser l'\'equiprobabilit\'e ?}]
L'hypoth\`ese d'\'equiprobabilit\'e n'est l\'egitime que si la sym\'etrie du
probl\`eme la justifie (d\'e \'equilibr\'e, tirage au sort\ldots). Ne l'invoquez
jamais sans v\'erification~!
\end{intuition}

\begin{exemple}[Probl\`eme des anniversaires]
Dans un groupe de $n$ personnes, quelle est la probabilit\'e qu'au moins deux
personnes aient le m\^eme anniversaire~? On suppose 365 jours \'equiprobables.

L'univers est $\Omega = \{1, \ldots, 365\}^n$ avec $\abs{\Omega} = 365^n$.
L'\'ev\'enement compl\'ementaire \og tous les anniversaires sont distincts \fg a
pour cardinal $365 \times 364 \times \cdots \times (365 - n + 1)$. Donc~:
\[
    \PP(\text{au moins une co\"incidence})
    = 1 - \frac{365!}{(365-n)!\, 365^n}.
\]
Pour $n = 23$, on obtient $\PP \approx 0{,}507$~: il y a plus d'une chance sur deux~!
\end{exemple}

\subsection{Probabilit\'e d\'efinie par une densit\'e discr\`ete}

\begin{proposition}
Si $\Omega$ est d\'enombrable et si $(p_\omega)_{\omega \in \Omega}$ est une famille
de r\'eels tels que $p_\omega \geq 0$ pour tout $\omega$ et
$\sum_{\omega \in \Omega} p_\omega = 1$, alors~:
\[
    \PP(A) = \sum_{\omega \in A} p_\omega, \quad A \subseteq \Omega,
\]
d\'efinit une mesure de probabilit\'e sur $(\Omega, \mathcal{P}(\Omega))$.
\end{proposition}

\subsection{Construction sur $\R$~: densit\'e continue}

\begin{proposition}
Soit $f : \R \to [0, +\infty)$ une fonction int\'egrable telle que
$\int_{-\infty}^{+\infty} f(x)\, dx = 1$. Alors~:
\[
    \PP(A) = \int_A f(x)\, dx, \quad A \in \mathcal{B}(\R),
\]
d\'efinit une mesure de probabilit\'e sur $(\R, \mathcal{B}(\R))$. La fonction $f$
est appel\'ee \textbf{densit\'e de probabilit\'e}.
\end{proposition}

% ------------------------------------------------------------
\section{Lemme de Borel--Cantelli (premi\`ere partie)}
% ------------------------------------------------------------

\begin{theoreme}[Premier lemme de Borel--Cantelli]
\label{thm:borel_cantelli_1}
Soit $(A_n)_{n \geq 1}$ une suite d'\'ev\'enements. Si
$\sum_{n=1}^{\infty} \PP(A_n) < \infty$, alors~:
\[
    \PP\!\left(\limsup_{n \to \infty} A_n\right) = 0,
\]
o\`u $\displaystyle\limsup_{n \to \infty} A_n
= \bigcap_{n=1}^{\infty} \bigcup_{k=n}^{\infty} A_k$ est l'\'ev\'enement
\og une infinit\'e de $A_n$ se r\'ealisent \fg.
\end{theoreme}

\begin{proof}
Posons $B_n = \bigcup_{k=n}^{\infty} A_k$. La suite $(B_n)$ est d\'ecroissante et
$\limsup A_n = \bigcap B_n$. Par continuit\'e d\'ecroissante~:
\[
    \PP(\limsup A_n) = \lim_{n \to \infty} \PP(B_n)
    \leq \lim_{n \to \infty} \sum_{k=n}^{\infty} \PP(A_k) = 0,
\]
car le reste d'une s\'erie convergente tend vers z\'ero.
\end{proof}

\begin{remarque}
Le \textbf{second} lemme de Borel--Cantelli (r\'eciproque partielle) n\'ecessite
l'ind\'ependance des \'ev\'enements et sera \'enonc\'e au chapitre~2.
\end{remarque}

% ------------------------------------------------------------
\section{Exercices}
% ------------------------------------------------------------

\begin{exercice}
Montrer que l'intersection de deux tribus est une tribu. L'union de deux tribus
est-elle toujours une tribu~? Donner un contre-exemple.
\end{exercice}

\begin{exercice}
Soit $\Omega = \{1,2,3,4,5,6\}$ (d\'e \'equilibr\'e). \'Ecrire la tribu
engendr\'ee par la partition $\bigl\{\{1,2\}, \{3,4\}, \{5,6\}\bigr\}$ et
d\'enombrer ses \'el\'ements.
\end{exercice}

\begin{exercice}
On lance un d\'e \'equilibr\'e deux fois. Calculer la probabilit\'e que~:
\begin{enumerate}[label=(\alph*)]
    \item la somme soit \'egale \`a~7\,;
    \item la somme soit sup\'erieure ou \'egale \`a~10\,;
    \item les deux faces soient identiques.
\end{enumerate}
\end{exercice}

\begin{exercice}
On tire trois cartes successivement et sans remise dans un jeu de 52~cartes.
Calculer la probabilit\'e d'obtenir exactement deux as.
\end{exercice}

\begin{exercice}
Soit $(A_n)$ une suite d\'ecroissante d'\'ev\'enements avec
$\PP(A_n) = 1/n$. Que vaut $\PP(\bigcap_{n=1}^{\infty} A_n)$~? Justifier.
\end{exercice}

\begin{exercice}[Lemme de Borel--Cantelli]
Soit $(A_n)$ une suite d'\'ev\'enements avec $\PP(A_n) = 1/n^2$.
Montrer que $\PP(\limsup A_n) = 0$. M\^eme question avec $\PP(A_n) = 1/n$~:
peut-on conclure~?
\end{exercice}

\begin{exercice}
Montrer la formule d'inclusion-exclusion (th\'eor\`eme~\ref{thm:inclusion_exclusion})
par r\'ecurrence sur~$n$.
\end{exercice}

\begin{exercice}
On dispose $n$ lettres dans $n$ enveloppes au hasard (une lettre par enveloppe).
Soit $D_n$ le nombre de lettres \og bien plac\'ees \fg (dans la bonne enveloppe).
En utilisant la formule d'inclusion-exclusion, montrer que~:
\[
    \PP(D_n = 0) = \sum_{k=0}^{n} \frac{(-1)^k}{k!}
    \xrightarrow[n \to \infty]{} e^{-1} \approx 0{,}368.
\]
\end{exercice}
